% Section 4: Mechanisms of useful context saturation
% Required packages: \usepackage{amsmath,booktabs,graphicx}
% Upload this file together with the four main figures under figures/.

\section{Mechanisms of useful context saturation}
\label{sec:context-mechanisms}

Section~\ref{sec:real-context} established that the context length that
minimizes forecast error is heterogeneous and often substantially shorter than
the available history.  An error--context curve alone, however, cannot explain
why accuracy saturates.  The distant past may be \emph{predictively redundant}
given recent observations, or the architecture may be unable to use distant
information even when it is necessary.  We distinguish these explanations with
three complementary tests: fixed-width temporal perturbations, oracle-verified
distant-dependency controls, and a decomposition of physical slicing versus
attention masking.

\subsection{Causal and mechanistic protocol}
\label{sec:mechanistic-protocol}

For this analysis, we define the $5\%$ sufficient context as
\begin{equation}
    L_{5\%}
    = \min\left\{
        L\in\mathcal{G}:
        E(L) \leq 1.05\min_{L'\in\mathcal{G}} E(L')
      \right\},
    \label{eq:sufficient-context}
\end{equation}
where $E(L)$ is mean forecast MAE and $\mathcal{G}$ is the model-supported
context grid.  Requested lengths are aligned to each model's patching and
context constraints, and all results below use the resulting effective length.
The $5\%$ tolerance prevents small fluctuations around a broad optimum from
being interpreted as substantively different windows.  Consequently,
$L_{5\%}$ is a performance threshold, not an assumed causal boundary.

The ordinary synthetic evaluation uses horizon 64 and seed 42.  Temporal
perturbations are available for 11 TSFMs, the complete planted-lag grid for 10,
and the masking--slicing decomposition for eight.  Each condition in the
completed planted-lag grid contains 128 instances; each masking--slicing cell
contains 256 series.  The long-lag extension contains 48 instances per
condition and is complete for nine models at $\alpha=1$, with an additional
$\alpha=0.5$ series for Toto.  We do not turn unsupported experiments or
incomplete cells into zero effects.

\subsection{Temporal sensitivity is strongly front-loaded}
\label{sec:temporal-perturbation-results}

For a temporal block $b$, input width $W$, and perturbation family $p$, we
measure
\begin{equation}
    A_p(b,W)
    = \mathbb{E}_i\!\left[
        \left|
        \ell\!\left(f(\widetilde{x}^{\,p}_{i,b}),y_i\right)
        - \ell\!\left(f(x_i),y_i\right)
        \right|
      \right].
    \label{eq:block-perturbation-effect}
\end{equation}
The intervention replaces only block $b$ while preserving the total input
width, positions, and mask.  We use block-mean replacement, within-block
permutation, a locally matched replacement, and Gaussian noise at three
severities.  The four families trade off intervention strength against
distribution shift; agreement across them is more informative than any single
corruption.

Figure~\ref{fig:perturbation-recency} compares block-mean profiles at each
model's largest completed input.  The dominant pattern is a steep decline in
sensitivity with lookback.  The three Chronos-2 variants, FlowState, Moirai,
PatchTST-FM, TimesFM, TiRex, Toto, and Chronos-Bolt all assign much smaller
effects to distant blocks than to the most recent block.  Sundial is a notable
exception: its profile falls initially but retains a broad plateau over older
blocks.  Useful context is therefore generally concentrated near the forecast
origin, but the concentration is architecture dependent.

\begin{figure*}[t]
    \centering
    \includegraphics[width=0.95\textwidth]
        {figures/fig1_perturbation_recency_all_models.png}
    \caption{Temporal block sensitivity across 11 TSFMs on the ordinary
    synthetic pool.  For each model, we show the block-mean intervention at
    its largest completed effective input $W$ and divide the mean absolute
    loss change by the largest value in that model's profile.  Lookback zero
    is displayed at one on the logarithmic axis.  The normalization supports
    comparison of profile shape, not absolute sensitivity magnitude.}
    \label{fig:perturbation-recency}
\end{figure*}

Chronos2-Small provides a fully scoped worked example.  Its mean MAE falls from
$0.797$ at $W=32$ to a minimum of $0.312$ at $W=3{,}072$, while the first
window within $5\%$ of that minimum is $L_{5\%}=1{,}536$.  At $W=8{,}192$, the
mean absolute effect outside versus inside this boundary is $0.146$ for
block-mean replacement, $0.189$ for matched replacement, $0.222$ for
permutation, and $0.376$ for noise.  Thus, the older prefix is much less
influential, but it is not absolutely inert.  The data support practical
concentration rather than a sharp causal cutoff at $1{,}536$.

\subsection{Verified distant dependencies induce adaptive reach}
\label{sec:planted-lag-results}

Low distant sensitivity on ordinary data remains compatible with architectural
inaccessibility.  We therefore construct controlled processes of the form
\begin{equation}
    x_t
    = g\!\left(x_{t-1:t-r}\right)
      + \alpha h(x_{t-d}) + \varepsilon_t,
    \qquad
    r=8,
    \quad
    d\in\{32,64,128,256\},
    \quad
    \alpha\in\{0,0.5,1\}.
    \label{eq:planted-lag-generator}
\end{equation}
A held-out ridge oracle must show that lag-$d$ features improve prediction over
recent-only features before the condition is used to draw a model conclusion.
No condition in the completed ten-model grid was marked as a broken generator.

For each model, we summarize the eight oracle-verified cells with
\begin{equation}
    \rho_m
    = \operatorname{Spearman}\!\left(d,L_{5\%}\right),
    \qquad
    q_m
    = \frac{1}{8}\sum_c
      \mathbf{1}\!\left[L_{5\%,c}\geq d_c\right],
    \label{eq:adaptive-reach-statistics}
\end{equation}
where the sum covers dependency strengths $\alpha\geq0.5$.  All ten models
have positive and strong lag tracking: $\rho_m$ ranges from $0.730$ to $0.908$.
The estimated sufficient context reaches the planted lag in every strong cell
for the three Chronos-2 variants, in seven of eight cells for six further
models, and in five of eight for Sundial (Table~\ref{tab:adaptive-reach}).

\begin{figure*}[t]
    \centering
    \includegraphics[width=0.98\textwidth]
        {figures/visual_draft/fig2_sufficient_context_vs_lag_with_toto_placeholder.pdf}
    \caption{Adaptive context reach across both oracle-verified dependency
    strengths.  Panels A--B show the completed ten-model grid at lags
    32--256.  Panel C shows the completed nine-model long-lag cohort at
    $\alpha=1$ and lags 128--2,048 with 48 series per condition; Toto's
    $\alpha=0.5$ series is retained as the only completed weak long-lag run.
    The dashed diagonal marks exact boundary recovery, and red rings identify
    conditions where $L_{5\%}<d$.}
    \label{fig:adaptive-reach}
\end{figure*}

\paragraph{Long-lag extension.}
At $\alpha=1$, the nine-model follow-up reaches the planted lag in 28 of 36
conditions.  Chronos2-Small, Chronos2-Synth, PatchTST-FM, TimesFM, and Toto
reach all four lags; Chronos2-Base reaches three, FlowState and Sundial two,
and Moirai one.  Toto selects $L_{5\%}=384,1{,}024,1{,}536,$ and $2{,}560$ at
$d=128,512,1{,}024,$ and $2{,}048$, respectively.  Its additional
$\alpha=0.5$ series selects $192,192,128,$ and $192$; only the first condition
reaches the planted lag.  Chronos-Bolt and TiRex do not have completed
long-lag controls in this cohort and are not shown in panel C.

\begin{table*}[t]
    \centering
    \footnotesize
    \setlength{\tabcolsep}{4pt}
    \caption{Lag tracking over the eight oracle-verified cells with
    $\alpha\geq0.5$.  ``Reach'' is $q_m$ in
    Equation~\ref{eq:adaptive-reach-statistics}; ``null range'' is the range of
    $L_{5\%}$ selected in the four $\alpha=0$ cells.}
    \label{tab:adaptive-reach}
    \begin{tabular}{lrrr}
        \toprule
        Model & $\rho_m$ & Reach & Null range \\
        \midrule
        Chronos2-Base      & 0.889 & 1.000 & 192--256 \\
        Chronos2-Small     & 0.889 & 1.000 & 128--256 \\
        Chronos2-Synth     & 0.908 & 1.000 & 128--256 \\
        ChronosBolt-Base   & 0.800 & 0.875 &  96--192 \\
        FlowState-R1       & 0.765 & 0.875 &  64--128 \\
        Moirai2-Small      & 0.800 & 0.875 &  96--192 \\
        PatchTST-FM-R1     & 0.765 & 0.875 &  48--192 \\
        Sundial-Base-128M  & 0.730 & 0.625 &  48--96  \\
        TiRex2             & 0.745 & 0.875 & 128--192 \\
        TimesFM2.5-200M    & 0.830 & 0.875 &  64--192 \\
        \bottomrule
    \end{tabular}
\end{table*}

The null cells expose an equally important limitation.  Although their distant
coefficient is zero, they select $L_{5\%}$ between 48 and 256 across models.
Local-only controls with true order eight select between 32 and 384.  Hence the
operational threshold tracks a real increase in dependency reach but does not
precisely recover the causal boundary.  This falsification test favors adaptive
context use over a fixed architectural limit, while showing that
$L_{5\%}$ is conservative and requires calibration.

\subsection{Masking effects are partly preprocessing effects}
\label{sec:masking-decomposition-results}

Physical slicing changes input width, positions, padding, and potentially
normalization statistics at the same time.  Conversely, hiding an old prefix
with an attention mask preserves width and position but may still allow the
hidden values to affect preprocessing.  For each visible suffix $L$ inside a
fixed full width $W$, we therefore compare
\begin{align}
    E_{\mathrm{slice}}(L)
        &: \text{physically supply only the last }L\text{ values}, \nonumber\\
    E_{\mathrm{fullstat}}(L)
        &: \text{mask the true old prefix at width }W, \nonumber\\
    E_{\mathrm{tailstat}}(L)
        &: \text{mask a prefix matched to the tail's mean and scale}.
    \label{eq:restriction-conditions}
\end{align}
Their differences define the operational decomposition
\begin{equation}
    \Delta_{\mathrm{stat}}(L)
      = E_{\mathrm{fullstat}}(L)-E_{\mathrm{tailstat}}(L),
    \qquad
    \Delta_{\mathrm{resid}}(L)
      = E_{\mathrm{tailstat}}(L)-E_{\mathrm{slice}}(L).
    \label{eq:restriction-decomposition}
\end{equation}
$\Delta_{\mathrm{stat}}$ exactly isolates normalization only for a global
affine mean--standard-deviation transform; for other pipelines it is a proxy
for normalization and related prefix-dependent preprocessing.

Figure~\ref{fig:restriction-decomposition} displays the three conditions
directly for every model with a completed decomposition.  Full-history-stat
masking can produce large and non-monotonic errors at short suffixes.  Matching
the hidden prefix to the visible tail removes almost all of this separation for
the Chronos variants, Chronos-Bolt, and Moirai, but not universally.

\begin{figure*}[t]
    \centering
    \includegraphics[width=\textwidth]
        {figures/visual_draft/fig3_slicing_vs_masking_all_models.pdf}
    \caption{Direct slicing--masking comparison across eight TSFMs.  Each
    panel has its own MAE scale.  Tail-matched masking is shown by teal dashed
    lines with open markers; where it lies directly over the physical-slice
    curve, the preprocessing-matched gap is visually negligible.  Material
    residual separations remain for PatchTST-FM, TimesFM, and Sundial.}
    \label{fig:restriction-decomposition}
\end{figure*}

For Chronos2-Small, mean $|\Delta\mathrm{MAE}|$ falls from $0.03115$ under
full-history-stat masking to $0.00105$ under tail matching, a $29.7\times$
difference.  Across eight models, tail matching reduces the mean absolute
masking--slicing gap by $48.2\%$ to approximately $100\%$, with a median
reduction of $97.2\%$.  The residual is not negligible everywhere: TimesFM
retains $51.8\%$ of the original gap, PatchTST-FM $34.5\%$, and Sundial
$17.3\%$; each of the other five models retains at most $3.4\%$.

\begin{figure*}[t]
    \centering
    \includegraphics[width=\textwidth]
        {figures/visual_draft/fig4_masking_gap_heatmap_all_models.pdf}
    \caption{Mean absolute masking--slicing MAE gap before and after tail-stat
    matching.  The logarithmic color scale makes broad preprocessing effects
    and smaller residuals simultaneously visible.  Pale cells denote suffixes
    not evaluated because the model's full context is shorter.}
    \label{fig:masking-gap-heatmap}
\end{figure*}

Apparent context use under fixed-width masking can therefore be dominated by
prefix-dependent preprocessing, but width, position, padding, or other
architecture-specific effects remain important for some models.  Masking is
not a neutral substitute for shortening the input: mechanistic claims should
report both full-stat and tail-matched conditions and treat the remaining gap
as architecture specific rather than universal.

\subsection{Implications for context selection}
\label{sec:mechanistic-implications}

Taken together, the experiments support an adaptive rather than fixed context
mechanism.  On ordinary series, forecast sensitivity is predominantly recent;
when a verified distant dependency becomes necessary, effective context grows
toward it across diverse architectures.  However, the operational
$L_{5\%}$ threshold overestimates causal span in null controls, and naive
masking can mistake preprocessing-statistic changes for information use.
Accordingly, the learned selector in the next section predicts performance over
a discrete context grid rather than a literal last-used timestep.  Its value is
to choose an accuracy--compute operating point, not to claim exact recovery of
the data-generating dependency length.
